\section{Tối ưu bằng QAT}

\subsection{Thiết lập quantization}

QAT được hiện thực bằng thư viện \texttt{pytorch-quantization} của NVIDIA, là
đường chính thức để tạo mô hình INT8 có thông tin lượng tử tường minh cho
TensorRT \cite{jacob2018}. Hàm \texttt{quant\_modules.initialize()} thay mọi
\texttt{nn.Conv2d} bằng \texttt{QuantConv2d} trước khi dựng mô hình. Theo mặc định
của thư viện, lượng tử hóa là \textbf{đối xứng}: trọng số lượng tử \textbf{theo
từng kênh đầu ra} (per-channel), activation lượng tử \textbf{theo tensor}
(per-tensor). Với các phép cộng dư và ghép kênh trong các khối
\texttt{C3k2}, \texttt{C2PSA} đã cắt tỉa, đồ án viết các lớp lượng tử riêng
(\texttt{quant\_ops\_pruned\_yolo26.py}) chèn \texttt{TensorQuantizer} cho từng
đầu vào, để TensorRT có thể hợp nhất phép toán ở INT8 thay vì chuyển đổi qua lại
giữa INT8 và FP16. Các lượng tử của những nhánh không cần thiết được vô hiệu hóa
để giữ độ chính xác.

\subsection{Hiệu chỉnh}

Trước khi huấn luyện, thang lượng tử (\texttt{amax}) của activation được hiệu
chỉnh trên 256 batch dữ liệu huấn luyện bằng phương pháp \textbf{entropy}: với
mỗi tensor, thu histogram giá trị và chọn ngưỡng cắt cực tiểu hóa KL divergence
giữa phân bố gốc và phân bố sau lượng tử. Phương pháp này bỏ qua các giá trị
ngoại lai hiếm gặp, cho độ phân giải lượng tử tốt hơn so với lấy giá trị lớn nhất.

\subsection{Fake quantization và huấn luyện}

Sau hiệu chỉnh, các nút fake-quantization mô phỏng phép làm tròn và cắt về INT8
trong lan truyền thuận, còn gradient đi qua bằng Straight-Through Estimator. Mô
hình được fine-tune nhẹ 20 epoch, batch 8, optimizer SGD momentum 0,9; learning
rate đầu vào $10^{-3}$ được chia 100 bên trong trainer (thành $10^{-5}$) vì QAT
chỉ cần điều chỉnh nhỏ quanh điểm đã hội tụ. Tùy chọn hiệu chỉnh lại định kỳ
(\texttt{recalib\_every}) mặc định tắt. QAT chỉ được chạy trên checkpoint \emph{đã
phục hồi} sau cắt tỉa, không chạy trên checkpoint vừa cắt còn yếu.

\subsection{Export INT8}

Checkpoint QAT được export sang ONNX, trong đó mỗi tensor lượng tử được biểu diễn
bằng một cặp nút Q/DQ (\texttt{QuantizeLinear} và \texttt{DequantizeLinear})
mang thang lượng tử đã học. TensorRT đọc các nút Q/DQ này để build engine INT8 trực tiếp, không cần bước
hiệu chỉnh PTQ riêng. Thư viện \texttt{pytorch-quantization} yêu cầu CUDA nên giai
đoạn này chạy trên môi trường Linux có GPU (Google Colab), không chạy được trên
máy Windows chỉ có CPU.

\subsection{Đánh giá mô hình QAT}

Mô hình QAT được đánh giá ở ba điểm: checkpoint fake-quant trong PyTorch, ONNX
sau export và TensorRT engine trên thiết bị. Chênh lệch giữa các điểm cho biết sai
số đến từ lượng tử, từ export hay từ backend.
